\begin{abstractzh}
本篇論文主要的目的是透過獲取騎乘者的生理訊息，與當下的環境資訊，選擇一個合適的檔位讓騎乘者可以順利地騎乘，不必顧慮檔位的切換問題。腳踏車平台我們使用先前建立的智慧自行車系統作為原型，該系統是一個可以即時的收集騎乘者生理狀態，與當下環境狀態的一個平台。我們選擇了幾種不同的決策方式對收集到的資料進行預測，透過收集到的感測器資訊對當下適合的檔位進行預測分析。由於有使用深度學習方面的方法，需要一定數量的資料用於進行訓練。因此將腳踏車系統稍加調整後，利用該系統收集騎乘者當下的多個感測器訊息，還有當下的檔位狀態，並且透過實際的騎乘收集一定數量的訓練資料。使用這些資料對我們所建立的機器學習方法進行訓練，最後透過一部分不包含在訓練資料內的資料對訓練後的結果進行測試。另外我們建立了一個網頁平台可以用於管理收集到的資料，該平台可以用於資料的上傳、查詢，還有簡單的分析結果可供查看。
\bigbreak
\noindent \textbf{關鍵字：}{\, \makeatletter \@keywordszh \makeatother}
\end{abstractzh}

\begin{abstracten}
The main purpose of this thesis is to obtain the physiological information of the rider, the current environmental information, and select a suitable gear so that the rider can ride smoothly without worrying about the problem of gear shifting. We use the previously established smart bicycle system as a prototype to collect the physiological state of the rider and the current environmental state in real time. 
Several different decision-making methods are used with the collected sensor information to predict and analyze the current suitable gear. Due to the use of deep learning methods, a certain amount of data is required for training. We use the previously established bicycle system collect the physiological information and environmental data of the rider. After slightly adjusting the system, we use the system to collect the current multiple sensor information of the rider, as well as the current gear status. 
These data are used to train the machine learning method we have established, and finally methods are used for testing. In addition, we have established a web platform that can be used to manage the collected data. The platform can be used for data upload, query, and simple analysis for viewing.


\bigbreak
\noindent \textbf{Keywords:}{\, \makeatletter \@keywordsen \makeatother}
\end{abstracten}

\begin{comment}
\category{I2.10}{Computing Methodologies}{Artificial Intelligence --
Vision and Scene Understanding} \category{H5.3}{Information
Systems}{Information Interfaces and Presentation (HCI) -- Web-based
Interaction.}

\terms{Design, Human factors, Performance.}

\keywords{Region of interest, Visual attention model, Web-based
games, Benchmarks.}
\end{comment}
